\clearpage
\section{Futurist: where the subject is going}
\label{sec:future}

The recent direction of environmental QC is not the replacement of simple
checks by one complicated model. It is the addition of better references
inside the same basic architecture.

First, QC is becoming software infrastructure. Tools such as PyHydroQC and
SaQC package detection, correction or configuration into reusable systems
\citep{jones2022,schmidt2023}. This matters because the configuration can
be versioned and inspected. A threshold is no longer only a number hidden in
code; it can become an auditable part of the data-quality policy.

Second, model-based anomaly detection is expanding quickly. Reviews such as
\citet{blazquez2021} describe deep learning and other modern time-series
methods. Their advantage is that they can learn more complicated versions of
normal behaviour than a fixed range or step limit. Their operational
challenge is explainability, transfer between sites, and the shortage of
well-labelled fault data.

Third, references are becoming more empirical. A spatial check can choose
neighbours based on measured agreement rather than distance alone. Seasonal
bands can be learned from validated history. Forecasting models can adapt to
a site's own behaviour. These changes improve the reference for ``normal'',
which is the part of the QC problem that has changed most across the five
readings.

I therefore expect future systems to remain hybrid: simple physical rules for
clear failures, robust and model-based methods for subtler problems, graded
flags, and human review for ambiguous cases. Better machine learning may
improve one member of the battery without removing the need for the battery.

% =====================================================================
% 6. Critic
% =====================================================================
